\chapter{Metodología}

\section{Caso de uso industrial}

El caso de uso seleccionado es la detección de anomalías en rodamientos de
motores eléctricos mediante señales de vibración. La elección permite trabajar
con datos industriales públicos, etiquetas interpretables y un flujo temporal
adecuado para ventanas, features y tensores.

\section{Pipeline de datos}

El flujo metodológico para el MVP queda definido como:

\begin{center}
\begin{minipage}{0.85\textwidth}
\begin{verbatim}
datos .mat crudos
-> manifiesto del dataset
-> extracción de señales
-> normalización de frecuencia de muestreo
-> perfilado estadístico
-> limpieza determinista
-> ventanas temporales
-> features y tensores
-> particiones train/validation/test
-> modelos de anomalías
-> evaluación
-> informe técnico
\end{verbatim}
\end{minipage}
\end{center}

\section{Generalización multi-dataset}

Tras consolidar el MVP sobre CWRU, la siguiente extensión metodológica consiste
en separar explícitamente tres niveles: el dataset crudo, el adaptador
determinista de dataset y el pipeline común. Esta separación evita que los
ejecutores dependan de nombres de fichero o estructuras particulares de un
benchmark concreto.

La generalización se plantea mediante un descriptor común de dataset, un
manifiesto común con una fila por unidad cruda procesable y un registro de
adaptadores. CWRU se mantendrá como benchmark de regresión, mientras que NASA
IMS se contempla como primer candidato para validar el sistema con series más
largas, múltiples canales y condiciones de degradación menos triviales.

La inspección local de NASA IMS muestra que el paquete se distribuye como un
contenedor anidado (`zip`, `7z` y `rar`) y que la aplicación deberá distinguir
entre detección de formato, selección de adaptador y generación efectiva de
manifiesto. En una futura interfaz, el usuario debería ver una previsualización
del descriptor del dataset, los formatos detectados, las limitaciones de
extracción y si el manifiesto puede generarse directamente o requiere preparar
una carpeta preextraída.

Como primer paso ejecutable, el manifiesto NASA IMS se restringe a carpetas ya
preextraídas. Cada snapshot temporal se representa como una fila común con
etiqueta inicial desconocida, fallo final conocido a nivel de set y metadatos
de frecuencia, canales y variante local. Esta decisión permite validar la
generalización del pipeline sin incorporar todavía dependencias de extracción
ni una política de etiquetado temporal.

Los agentes no interpretarán directamente carpetas crudas. Recibirán
descriptores, resúmenes de manifiesto y diagnósticos de calidad generados por
código determinista. De este modo pueden aumentar su capacidad de decisión
--por ejemplo, recomendando canal, ventana, familia de features o modelo-- sin
ejecutar código arbitrario ni saltarse los contratos del sistema.

\section{Manifiesto del dataset}

El primer artefacto metodológico será un manifiesto con una fila por fichero
crudo. Este manifiesto conectará los `.mat` originales con el resto del
pipeline. Las columnas previstas son:

\begin{longtable}{ll}
\toprule
Campo & Descripción \\
\midrule
file\_id & identificador del fichero original \\
dataset & nombre del dataset \\
source\_path & ruta del fichero crudo \\
label & clase `normal' o `fault' \\
fault\_type & tipo de fallo si existe \\
fault\_diameter\_inch & diámetro del defecto \\
load\_hp & carga del motor \\
rpm & régimen de giro \\
sensor\_channel & canal de sensor usado \\
source\_sample\_rate\_hz & frecuencia de muestreo original \\
target\_sample\_rate\_hz & frecuencia objetivo del MVP \\
source\_format & formato original \\
notes & observaciones \\
\bottomrule
\end{longtable}

\section{Extracción y perfilado}

El canal principal inicial será `DE\_time`, asociado al acelerómetro del lado
drive-end. El perfilado generará resúmenes estadísticos por fichero y por
etiqueta, incluyendo número de muestras, canales detectados, valores no finitos,
media, desviación típica, RMS, curtosis, factor de cresta y energía.

\section{Estructuración temporal}

La frecuencia objetivo inicial será 12 kHz. Las ventanas candidatas serán de
2048 o 4096 muestras con solapamiento del 50\%. Esta decisión permite comparar
features tabulares con modelos clásicos y conservar ventanas crudas para
autoencoders posteriores.

\section{Particiones}

Para evitar fuga de información entre ventanas, no se usará un split aleatorio
ingenuo. La estrategia inicial será:

\begin{itemize}
  \item entrenamiento con ventanas normales;
  \item validación con ventanas normales retenidas;
  \item test con ventanas normales retenidas y ventanas con fallo.
\end{itemize}

\section{Métricas}

Las métricas iniciales serán precision, recall, F1-score, ROC-AUC, PR-AUC, tasa
de falsos positivos y matriz de confusión. Desde el punto de vista industrial,
el criterio principal será mantener alta sensibilidad ante fallos sin disparar
en exceso los falsos positivos.
